\section{若尔当--赫尔德定理}

在可解群的讨论中，我们利用相邻商群均为交换群的递降子群列描述群的结构。对于一般有限群，也可用这样的子群列逐步分解，只是相邻商群不必交换。如果一直细分到无法再插入新的项，便得到以单群为因子的群列。下面研究这种分解在多大程度上由原群决定。

\begin{definition}
设有限递降子群列
\[
G=G_0\geq G_1\geq\cdots\geq G_r=\{e\}
\]
满足$G_i\trianglelefteq G_{i-1}$，则称其为$G$的一个\textbf{次正规子群列}。各商群$G_{i-1}/G_i$组成的群组称为其\textbf{因子群组}。这里的“次”表示只要求每一项在前一项中正规，不要求它在整个$G$中正规。

允许群列中有重复项，重复项对应平凡因子群；非平凡因子群的个数称为该群列的\textbf{长度}。若群列无重复项，且各因子群均为单群，则称它为\textbf{合成群列}。
\end{definition}

若$G$本身为单群，则$G>\{e\}$就是唯一的无重复项合成群列。一般群的次正规子群列并不唯一；我们先证明有限群至少存在一条合成群列。

\begin{proposition}
每个有限群都有合成群列；无限群则未必如此。
\end{proposition}
\begin{proof}
平凡群可取长度为0的群列。设$G\ne\{e\}$。在所有无重复项的次正规子群列中，选择长度最大的一条。由于每次严格下降都使子群阶变小，这些长度有有限上界，故这样的选择存在。

若其中某个因子$G_i/G_{i+1}$不是单群，则它有非平凡真正规子群$\overline H$。由子群对应定理，完全反像$H$满足
\[
G_{i+1}<H<G_i,\qquad H\trianglelefteq G_i,
\qquad H/G_{i+1}=\overline H.
\]
而$G_{i+1}\trianglelefteq G_i$也保证$G_{i+1}\trianglelefteq H$，所以可以把$H$插入群列，形成一条更长的无重复项次正规子群列，矛盾。因此每个因子均单，所选群列为合成群列。
\end{proof}

这一论证还说明，合成群列恰是不能在相邻两项之间真正插入新子群的次正规子群列。有限性在这里保证插入过程最终结束。

\begin{example}
整数加法群$\Z$没有有限的合成群列。事实上，若有这样一条群列，其最后一个非平凡项必须是$\Z$的单子群。但$\Z$的每个非零子群都是$m\Z$，同构于$\Z$；它又含非平凡真子群$2m\Z$，所以不可能是单群。
\end{example}

第四节已证明，有限可解群的合成因子都是素数阶循环群，且这些素数计入重数后由群阶确定。对于一般有限群，也有相应的唯一性，只是因子可能为非交换单群。

\begin{theorem}\label{p2:jordan-holder}
\textbf{若尔当--赫尔德定理：}有限群$G$的任意两条无重复项合成群列长度相等，其因子群组在不计次序、保留重复次数的意义下对应同构。
\end{theorem}
\begin{proof}
设两条合成群列为
\begin{align*}
G&=G_0\triangleright G_1\triangleright\cdots\triangleright G_r=\{e\},\\
G&=H_0\triangleright H_1\triangleright\cdots\triangleright H_s=\{e\}.
\end{align*}
对第一条的长度$r$归纳。平凡群的情形成立；当$r=1$时，$G$为单群，它没有可作第一步的非平凡真正规子群，所以第二条也只能为$G>\{e\}$，结论成立。

设$r>1$，并假设定理对具有长度小于$r$的合成群列的有限群都成立。

\textbf{若$G_1=H_1$：}删去两列的首项$G$，便得到同一群$G_1$的两条合成群列，其中第一条长度为$r-1$。归纳假设给出$r-1=s-1$，且下面的因子群组对应同构。补回共同的第一个因子$G/G_1$，即得结论。

\textbf{若$G_1\ne H_1$：}由于$G/G_1$和$G/H_1$均为单群，$G_1,H_1$都是$G$的极大真正规子群。它们不能彼此包含，否则极大性将迫使二者相等。乘积$G_1H_1$是包含两者的正规子群，严格大于$G_1$，所以
\[
G_1H_1=G.
\]
令$N_2=G_1\cap H_1$。由第二同构定理，
\[
G_1/N_2\cong G/H_1,\qquad
H_1/N_2\cong G/G_1.
\]
因此这两个商群均为单群。

取$N_2$的一条合成群列$N_2>N_3>\cdots>N_t=\{e\}$，若$N_2=\{e\}$，则取$t=2$。接成两条新合成群列：
\begin{align*}
G&\triangleright G_1\triangleright N_2\triangleright\cdots\triangleright N_t=\{e\},\\
G&\triangleright H_1\triangleright N_2\triangleright\cdots\triangleright N_t=\{e\}.
\end{align*}
它们都是合成群列，因为最上面两个商群分别由已知的单群给出，下面沿用$N_2$的合成群列。

先比较原第一条群列与第一条新群列。二者从$G_1$开始都是$G_1$的合成群列，原群列的这段长度为$r-1$，可用归纳假设。因此第一条新群列也有长度$r$，即$t=r$，且它的因子群组与原第一条对应同构。

第二条新群列与第一条新群列使用相同的下半段，故长度也为$r$。再比较它从$H_1$开始的群列与原第二条从$H_1$开始的群列，前者长度为$r-1$，又可用归纳假设，得到$s=r$，并且原第二条与第二条新群列的因子群组对应同构。

最后，两条新群列从$N_2$以下的因子相同，而最上面两个因子分别为
\[
G/G_1,\ G_1/N_2
\quad\text{及}\quad
G/H_1,\ H_1/N_2.
\]
由前面的两个同构，它们只是互换次序。因此两条新群列的因子群组对应同构，进而原来的两条群列也如此，归纳完成。
\end{proof}

\begin{example}
在$S_4$中可以取
\[
S_4\triangleright A_4\triangleright V_4\triangleright C_2\triangleright\{e\},
\]
其中$C_2$是$V_4$的任意一个二阶子群。各因子依次同构于$C_2,C_3,C_2,C_2$。这里$C_2$只须在$V_4$中正规，不要求在$S_4$中正规。

$C_4$与$V_4$都有两个同构于$C_2$的合成因子，但前者有四阶元素，后者没有，所以二者不同构。这说明合成因子群组并不能反过来唯一确定群本身。
\end{example}

更一般地，给定群$H,N$，寻找所有具有正规子群同构于$N$、相应商群同构于$H$的群$G$，即研究
\[
G/N\cong H,
\]
所允许的群结构，称为\textbf{群扩张问题}。合成群列把群分解为单群因子，但从这些因子重新构造原群，还需要知道它们如何结合。
